% Part: first-order-logic
% Chapter: proof-systems
% Section: natural-deduction

\documentclass[../../../include/open-logic-section]{subfiles}

\begin{document}

\iftag{FOL}
      {\olfileid{fol}{prf}{ntd}}
      {\olfileid{pl}{prf}{ntd}}

\olsection{طبيعي استنتاج}

طبيعي استنتاج د !!a{derivation} داسې نظام دے چې موخه ئې د رښتيني
استدلال، په تېره بيا د رياضي پوهانو د منظم استدلال، انځورول دي۔
رښتينے استدلال د يو شمېر ``طبيعي'' بڼو له مخې مخکښې ځي۔ مثلاً د
حالتونو ثبوت موږ ته اجازه راکوي چې د يوه فصلي فرض پر بنسټ پايله
ثابته کړو، په دې چې وښيو پايله د فصل له هرې برخې څخه راځي۔ نامستقيم
ثبوت اجازه راکوي چې پايله په دې سره ثابته کړو چې نفي ئې تناقض ته
رسېږي۔ شرطي ثبوت د ``که \dots نو \dots'' شرطي ادعا په دې سره ثابتوي
چې تالي له مقدم څخه راځي۔ طبيعي استنتاج د دغو طبيعي استنتاجونو د
ځينو صوري کول دي۔ هر منطقي نښلوونکے او هر سور دوه قاعدې لري، يوه د
معرفي کولو او بله د لرې کولو قاعده، او هره يوه ئې د داسې طبيعي
استنتاج له يوې بڼې سره سمه ده۔ مثلاً $\Intro{\lif}$ له شرطي ثبوت
سره او $\Elim{\lor}$ د حالتونو له ثبوت سره سمون لري۔ يوه په تېره
ساده قاعده $\Elim{\land}$ ده، چې اجازه راکوي له $!A \land !B$ څخه
~$!A$ يا $!B$ استنتاج کړو۔

يو خاصيت چې طبيعي استنتاج د !!{derivation} له نورو نظامونو څخه
بېلوي، د فرضونو کارول دي۔ په طبيعي استنتاج کښې !!^a{derivation} د
!!{formula}s يوه ونه ده۔ يو !!{formula} د !!{formula}s د ونې په جرړه
کښې وي، او د ونې ``پاڼې'' هغه !!{formula}s وي چې پايله ترې راايستل
کېږي۔ په طبيعي استنتاج کښې ځينې پاڼه-!!{formula}s د !!{derivation}
دننه ونډه لري، خو تر هغې ``کارول شوي وي'' چې !!{derivation} پايلې ته
رسېږي۔ دا په رښتيني استدلال کښې د هغو فرضيو د معرفي کولو له دود سره
سمون لري چې يوازې لږ وخت نافذې پاتې کېږي۔ مثلاً د حالتونو په ثبوت
کښې د فصل د هرې برخې رښتياوالے فرضوو؛ په شرطي ثبوت کښې د مقدم
رښتياوالے فرضوو؛ او په نامستقيم ثبوت کښې د پايلې د نفي رښتياوالے
فرضوو۔ د فرضي فرضونو د معرفي کولو او بيا د يوه منځني ګام د ثابتولو
دپاره د هغو ختمولو دا طريقه د طبيعي استنتاج څرګنده نښه ده۔ د طبيعي
استنتاج د يوه !!{derivation} په پاڼو کښې فارمولونه فرضونه بلل کېږي،
او د استنتاج ځينې قاعدې کېدے شي هغه ``!!{discharge}'' کړي۔ مثلاً که
مو له ځينو فرضونو څخه د~$!B$ کوم !!a{derivation} درلود چې~$!A$ هم
پکې شامل و، نو د $\Intro{\lif}$ قاعده اجازه راکوي چې~$!A \lif !B$
استنتاج کړو او د~$!A$ بڼې هر فرض ختم کړو۔ (د دې ثبتولو
دپاره چې کوم فرضونه په کومو استنتاجونو کښې ختمېږي، استنتاج او هغه
فرضونه چې ختموي ئې په يوه شمېره نښه کوو۔) هغه فرضونه چې د
!!{derivation} په پاے کښې !!{undischarged} پاتې کېږي، په ګډه د پايلې
د رښتياوالي دپاره بس دي؛ نو !!a{derivation} ثابتوي چې د هغه
!!{undischarged} فرضونه پايلې ته استلزام ورکوي۔

د طبيعي استنتاج پر بنسټ اړيکه $\Gamma \Proves !A$ هله صدق کوي چې
داسې !!a{derivation} موجود وي چې $!A$ د ونې وروستۍ !!{sentence} وي،
او هره !!{undischarged} پاڼه په~$\Gamma$ کښې وي۔ $!A$ په طبيعي
استنتاج کښې هله قضيه ده چې داسې !!a{derivation} موجود وي چې $!A$
وروستۍ !!{sentence} وي او ټول فرضونه !!{discharged} وي۔ مثلاً دا
!!a{derivation} ښيي چې $\Proves (!A \land !B) \lif !A$:
\begin{prooftree}
  \AxiomC{$\Discharge{!A \land !B}{1}$}
  \RightLabel{\Elim{\land}}
  \UnaryInfC{$!A$}
  \DischargeRule{\Intro{\lif}}{1}
  \UnaryInfC{$(!A \land !B) \lif !A$}
\end{prooftree}
نښه~$1$ څرګندوي چې فرض $!A \land !B$ د \Intro{\lif} په استنتاج کښې
!!{discharged} کېږي۔
  
يو سټ~$\Gamma$ هله ناسازګار دے چې په طبيعي استنتاج کښې
$\Gamma \Proves \lfalse$ وي۔ د \FalseInt{} قاعده دا يقيني کوي چې له
ناسازګار سټ څخه هر !!{sentence} !!{derive}d کېدے شي۔

د طبيعي استنتاج نظامونه ګيرهارډ ګېنتڅن او ستانيسواف ياشکوفسکي په
۱۹۳۰مو کلونو کښې جوړ کړل، او وروسته ډاګ پراوېڅ او فريډرېک فېچ لا
پراخ کړل۔ ځکه چې استنتاجونه ئې د ثبوت طبيعي لارې انځوروي، فلسفيان
ئې خوښوي۔ د فېچ جوړې شوې بڼې زياتره د منطق په پېژندنيزو درسي
کتابونو کښې کارول کېږي۔ د منطق په فلسفه کښې کله کله د طبيعي استنتاج
قاعدې د منطقي عملګرو معنٰی ورکوي، چې دې ته ``ثبوت-نظري معنايي
پوهنه'' ويل کېږي۔

\end{document}
